% Populärvetenskaplig sammanfattning på svenska
% Skriv denna text så att den är begriplig för en bred publik utan teknisk bakgrund

\section*{Bakgrund}
Östersjön är ett viktigt handelsområde med intensiv sjötrafik. Men det finns också fartyg som opererar utanför de normala regelverken—så kallade skuggflottan. Dessa fartyg kan användas för att kringgå internationella sanktioner, smuggla varor, eller bedriva andra illegala aktiviteter. De utgör ett hot mot sjöfartssäkerhet, miljön och internationell handel.

Skuggflottans fartyg försöker undvika upptäckt genom att manipulera sina spårningssystem, byta flagg ofta, eller stänga av sina identifieringssystem när de genomför misstänkta aktiviteter. Detta gör det svårt för myndigheter att övervaka och kontrollera dem.

\section*{Problem}
Traditionell övervakning av fartyg görs manuellt av analytiker som granskar rörelsemönster och letar efter misstänkt beteende. Men med tusentals fartyg i Östersjön varje dag är detta omöjligt att göra effektivt. Det behövs därför automatiserade system som kan hjälpa till att identifiera misstänkta fartyg.

\section*{Metod}
I detta examensarbete har jag utvecklat ett system som använder maskininlärning för att automatiskt upptäcka fartyg som uppvisar beteenden typiska för skuggflottan. Systemet analyserar data från AIS (Automatic Identification System), som är det spårningssystem som de flesta fartyg använder för att sända sin position.

Genom att studera rörelsemönster från kända skuggflottansfartyg har jag tränat en datormodell att känna igen karakteristiska beteenden. Modellen tittar på faktorer som:
\begin{itemize}
    \item Oregelbundna rutter och ovanliga destinationer
    \item Plötsliga förändringar i kurs eller hastighet
    \item Tidsperioder då spårningssystemet är avstängt
    \item Fartygets historik av flaggbyten
    \item Besök i hamnar eller områden kopplade till illegal verksamhet
\end{itemize}

\section*{Resultat}
[Fyll i dina faktiska resultat här när du har dem]

Systemet har testats på historisk data från Östersjön och visar lovande resultat i att identifiera misstänkta fartyg. Modellen kan [beskriv prestanda, t.ex. "upptäcka X\% av kända skuggflottansfartyg med en falsk alarm-frekvens på Y\%"].

\section*{Betydelse}
Detta arbete visar att maskininlärning kan vara ett värdefullt verktyg för maritim säkerhet. Ett automatiserat system kan hjälpa myndigheter att:
\begin{itemize}
    \item Snabbt identifiera misstänkta fartyg bland tusentals legitima fartyg
    \item Fokusera sina begränsade resurser på de fartyg som verkligen kräver granskning
    \item Upptäcka nya mönster i skuggflottans verksamhet
\end{itemize}

Systemet skulle kunna integreras i befintliga övervakningssystem och ge realtidsvarningar när misstänkta fartyg upptäcks.

\section*{Framtida utveckling}
För att göra systemet ännu bättre kan framtida arbete inkludera fler datakällor (t.ex. satellitbilder), mer avancerade maskininlärningsmodeller, och integration med andra övervakningssystem.
